\documentclass[12pt,]{article}
\usepackage{lmodern}
\usepackage{amssymb,amsmath}
\usepackage{ifxetex,ifluatex}
\usepackage{fixltx2e} % provides \textsubscript
\ifnum 0\ifxetex 1\fi\ifluatex 1\fi=0 % if pdftex
  \usepackage[T1]{fontenc}
  \usepackage[utf8]{inputenc}
\else % if luatex or xelatex
  \ifxetex
    \usepackage{mathspec}
  \else
    \usepackage{fontspec}
  \fi
  \defaultfontfeatures{Ligatures=TeX,Scale=MatchLowercase}
\fi
% use upquote if available, for straight quotes in verbatim environments
\IfFileExists{upquote.sty}{\usepackage{upquote}}{}
% use microtype if available
\IfFileExists{microtype.sty}{%
\usepackage{microtype}
\UseMicrotypeSet[protrusion]{basicmath} % disable protrusion for tt fonts
}{}
\usepackage[margin=1in]{geometry}
\usepackage{hyperref}
\hypersetup{unicode=true,
            pdftitle={Exceedingly Simple Sorting with Indices},
            pdfauthor={Jan de Leeuw},
            pdfborder={0 0 0},
            breaklinks=true}
\urlstyle{same}  % don't use monospace font for urls
\usepackage{color}
\usepackage{fancyvrb}
\newcommand{\VerbBar}{|}
\newcommand{\VERB}{\Verb[commandchars=\\\{\}]}
\DefineVerbatimEnvironment{Highlighting}{Verbatim}{commandchars=\\\{\}}
% Add ',fontsize=\small' for more characters per line
\usepackage{framed}
\definecolor{shadecolor}{RGB}{248,248,248}
\newenvironment{Shaded}{\begin{snugshade}}{\end{snugshade}}
\newcommand{\KeywordTok}[1]{\textcolor[rgb]{0.13,0.29,0.53}{\textbf{#1}}}
\newcommand{\DataTypeTok}[1]{\textcolor[rgb]{0.13,0.29,0.53}{#1}}
\newcommand{\DecValTok}[1]{\textcolor[rgb]{0.00,0.00,0.81}{#1}}
\newcommand{\BaseNTok}[1]{\textcolor[rgb]{0.00,0.00,0.81}{#1}}
\newcommand{\FloatTok}[1]{\textcolor[rgb]{0.00,0.00,0.81}{#1}}
\newcommand{\ConstantTok}[1]{\textcolor[rgb]{0.00,0.00,0.00}{#1}}
\newcommand{\CharTok}[1]{\textcolor[rgb]{0.31,0.60,0.02}{#1}}
\newcommand{\SpecialCharTok}[1]{\textcolor[rgb]{0.00,0.00,0.00}{#1}}
\newcommand{\StringTok}[1]{\textcolor[rgb]{0.31,0.60,0.02}{#1}}
\newcommand{\VerbatimStringTok}[1]{\textcolor[rgb]{0.31,0.60,0.02}{#1}}
\newcommand{\SpecialStringTok}[1]{\textcolor[rgb]{0.31,0.60,0.02}{#1}}
\newcommand{\ImportTok}[1]{#1}
\newcommand{\CommentTok}[1]{\textcolor[rgb]{0.56,0.35,0.01}{\textit{#1}}}
\newcommand{\DocumentationTok}[1]{\textcolor[rgb]{0.56,0.35,0.01}{\textbf{\textit{#1}}}}
\newcommand{\AnnotationTok}[1]{\textcolor[rgb]{0.56,0.35,0.01}{\textbf{\textit{#1}}}}
\newcommand{\CommentVarTok}[1]{\textcolor[rgb]{0.56,0.35,0.01}{\textbf{\textit{#1}}}}
\newcommand{\OtherTok}[1]{\textcolor[rgb]{0.56,0.35,0.01}{#1}}
\newcommand{\FunctionTok}[1]{\textcolor[rgb]{0.00,0.00,0.00}{#1}}
\newcommand{\VariableTok}[1]{\textcolor[rgb]{0.00,0.00,0.00}{#1}}
\newcommand{\ControlFlowTok}[1]{\textcolor[rgb]{0.13,0.29,0.53}{\textbf{#1}}}
\newcommand{\OperatorTok}[1]{\textcolor[rgb]{0.81,0.36,0.00}{\textbf{#1}}}
\newcommand{\BuiltInTok}[1]{#1}
\newcommand{\ExtensionTok}[1]{#1}
\newcommand{\PreprocessorTok}[1]{\textcolor[rgb]{0.56,0.35,0.01}{\textit{#1}}}
\newcommand{\AttributeTok}[1]{\textcolor[rgb]{0.77,0.63,0.00}{#1}}
\newcommand{\RegionMarkerTok}[1]{#1}
\newcommand{\InformationTok}[1]{\textcolor[rgb]{0.56,0.35,0.01}{\textbf{\textit{#1}}}}
\newcommand{\WarningTok}[1]{\textcolor[rgb]{0.56,0.35,0.01}{\textbf{\textit{#1}}}}
\newcommand{\AlertTok}[1]{\textcolor[rgb]{0.94,0.16,0.16}{#1}}
\newcommand{\ErrorTok}[1]{\textcolor[rgb]{0.64,0.00,0.00}{\textbf{#1}}}
\newcommand{\NormalTok}[1]{#1}
\usepackage{graphicx,grffile}
\makeatletter
\def\maxwidth{\ifdim\Gin@nat@width>\linewidth\linewidth\else\Gin@nat@width\fi}
\def\maxheight{\ifdim\Gin@nat@height>\textheight\textheight\else\Gin@nat@height\fi}
\makeatother
% Scale images if necessary, so that they will not overflow the page
% margins by default, and it is still possible to overwrite the defaults
% using explicit options in \includegraphics[width, height, ...]{}
\setkeys{Gin}{width=\maxwidth,height=\maxheight,keepaspectratio}
\IfFileExists{parskip.sty}{%
\usepackage{parskip}
}{% else
\setlength{\parindent}{0pt}
\setlength{\parskip}{6pt plus 2pt minus 1pt}
}
\setlength{\emergencystretch}{3em}  % prevent overfull lines
\providecommand{\tightlist}{%
  \setlength{\itemsep}{0pt}\setlength{\parskip}{0pt}}
\setcounter{secnumdepth}{5}
% Redefines (sub)paragraphs to behave more like sections
\ifx\paragraph\undefined\else
\let\oldparagraph\paragraph
\renewcommand{\paragraph}[1]{\oldparagraph{#1}\mbox{}}
\fi
\ifx\subparagraph\undefined\else
\let\oldsubparagraph\subparagraph
\renewcommand{\subparagraph}[1]{\oldsubparagraph{#1}\mbox{}}
\fi

%%% Use protect on footnotes to avoid problems with footnotes in titles
\let\rmarkdownfootnote\footnote%
\def\footnote{\protect\rmarkdownfootnote}

%%% Change title format to be more compact
\usepackage{titling}

% Create subtitle command for use in maketitle
\providecommand{\subtitle}[1]{
  \posttitle{
    \begin{center}\large#1\end{center}
    }
}

\setlength{\droptitle}{-2em}

  \title{Exceedingly Simple Sorting with Indices}
    \pretitle{\vspace{\droptitle}\centering\huge}
  \posttitle{\par}
    \author{Jan de Leeuw}
    \preauthor{\centering\large\emph}
  \postauthor{\par}
      \predate{\centering\large\emph}
  \postdate{\par}
    \date{Version 01, March 22, 2017}


\begin{document}
\maketitle
\begin{abstract}
We use the system \texttt{qsort} to write a routine that produces both
the sort an the order of a vector of doubles.
\end{abstract}

{
\setcounter{tocdepth}{3}
\tableofcontents
}
Note: This is a working paper which will be expanded/updated frequently.
All suggestions for improvement are welcome. The directory
\href{http://deleeuwpdx.net/pubfolders/mysort}{deleeuwpdx.net/pubfolders/mysort}
has a pdf version, the complete Rmd file with the code chunks, and the R
source code.

\section{Introduction}\label{introduction}

This may not be of much use to anyone else, but I'll make it available
anyway. For a larger project I needed something simple in C to sort a
vector and return both the sorted vector and the corresponding index
vector. Thus it combines \texttt{sort} and \texttt{order} in R. It
creates an array of structures in C, where each structure has a value
member and an index member. It then sorts the structures using the
system \texttt{qsort}, by comparison of the values. Finally it splits
the sorted array of structures into values and indices and returns those
to R in a two-element list.

\section{Example}\label{example}

\begin{Shaded}
\begin{Highlighting}[]
\KeywordTok{set.seed}\NormalTok{ (}\DecValTok{12345}\NormalTok{)}
\NormalTok{x <-}\StringTok{ }\KeywordTok{rnorm}\NormalTok{ (}\DecValTok{10}\NormalTok{)}
\KeywordTok{mySort}\NormalTok{ (x)}
\end{Highlighting}
\end{Shaded}

\begin{verbatim}
## $values
##  [1] -1.8179559677 -0.9193220025 -0.4534971735 -0.2841597439 -0.2761841052
##  [6] -0.1093033147  0.5855288178  0.6058874558  0.6300985511  0.7094660175
## 
## $indices
##  [1]  6 10  4  9  8  3  1  5  7  2
\end{verbatim}

as compared to

\begin{Shaded}
\begin{Highlighting}[]
\KeywordTok{list}\NormalTok{ (}\DataTypeTok{values =} \KeywordTok{sort}\NormalTok{ (x), }\DataTypeTok{indices =} \KeywordTok{order}\NormalTok{ (x))}
\end{Highlighting}
\end{Shaded}

\begin{verbatim}
## $values
##  [1] -1.8179559677 -0.9193220025 -0.4534971735 -0.2841597439 -0.2761841052
##  [6] -0.1093033147  0.5855288178  0.6058874558  0.6300985511  0.7094660175
## 
## $indices
##  [1]  6 10  4  9  8  3  1  5  7  2
\end{verbatim}

A quick comparison shows very little timing difference.

\begin{Shaded}
\begin{Highlighting}[]
\NormalTok{f <-}\StringTok{ }\ControlFlowTok{function}\NormalTok{ () \{}
\NormalTok{  x <-}\StringTok{ }\KeywordTok{rnorm}\NormalTok{(}\DecValTok{10000}\NormalTok{)}
\NormalTok{  h <-}\StringTok{ }\KeywordTok{mySort}\NormalTok{ (x)}
\NormalTok{\}}
\NormalTok{g <-}\StringTok{ }\ControlFlowTok{function}\NormalTok{ () \{}
\NormalTok{  x <-}\StringTok{ }\KeywordTok{rnorm}\NormalTok{ (}\DecValTok{10000}\NormalTok{)}
\NormalTok{  h <-}\StringTok{ }\KeywordTok{list}\NormalTok{ (}\DataTypeTok{values =} \KeywordTok{sort}\NormalTok{ (x), }\DataTypeTok{indices =} \KeywordTok{order}\NormalTok{ (x))}
\NormalTok{\}}
\KeywordTok{microbenchmark}\NormalTok{ (f, g, }\DataTypeTok{times =} \DecValTok{1000}\NormalTok{)}
\end{Highlighting}
\end{Shaded}

\begin{verbatim}
## Unit: nanoseconds
##  expr min lq   mean median uq  max neval cld
##     f  29 31 32.423     32 32  668  1000   a
##     g  29 31 39.067     32 32 7258  1000   a
\end{verbatim}

\section{C code}\label{c-code}

\begin{Shaded}
\begin{Highlighting}[]
\PreprocessorTok{#include }\ImportTok{<math.h>}
\PreprocessorTok{#include }\ImportTok{<stdlib.h>}

\DataTypeTok{int}\NormalTok{ myComp(}\DataTypeTok{const} \DataTypeTok{void}\NormalTok{ *, }\DataTypeTok{const} \DataTypeTok{void}\NormalTok{ *);}
\DataTypeTok{void}\NormalTok{ mySort(}\DataTypeTok{double}\NormalTok{ *, }\DataTypeTok{int}\NormalTok{ *, }\DataTypeTok{int}\NormalTok{ *);}

\KeywordTok{struct}\NormalTok{ couple \{}
    \DataTypeTok{double}\NormalTok{ value;}
    \DataTypeTok{int}\NormalTok{ index;}
\NormalTok{\};}

\DataTypeTok{int}\NormalTok{ myComp(}\DataTypeTok{const} \DataTypeTok{void}\NormalTok{ *px, }\DataTypeTok{const} \DataTypeTok{void}\NormalTok{ *py) \{}
    \DataTypeTok{double}\NormalTok{ x = ((}\KeywordTok{struct}\NormalTok{ couple *)px)->value;}
    \DataTypeTok{double}\NormalTok{ y = ((}\KeywordTok{struct}\NormalTok{ couple *)py)->value;}
    \ControlFlowTok{return}\NormalTok{ (}\DataTypeTok{int}\NormalTok{)copysign(}\FloatTok{1.0}\NormalTok{, x - y);}
\NormalTok{\}}

\DataTypeTok{void}\NormalTok{ mySort(}\DataTypeTok{double}\NormalTok{ *x, }\DataTypeTok{int}\NormalTok{ *k, }\DataTypeTok{int}\NormalTok{ *n) \{}
    \DataTypeTok{int}\NormalTok{ nn = *n;}
    \KeywordTok{struct}\NormalTok{ couple *xi =}
\NormalTok{        (}\KeywordTok{struct}\NormalTok{ couple *)calloc((}\DataTypeTok{size_t}\NormalTok{)nn, (}\DataTypeTok{size_t}\NormalTok{)}\KeywordTok{sizeof}\NormalTok{(}\KeywordTok{struct}\NormalTok{ couple));}
    \ControlFlowTok{for}\NormalTok{ (}\DataTypeTok{int}\NormalTok{ i = }\DecValTok{0}\NormalTok{; i < nn; i++) \{}
\NormalTok{        xi[i].value = x[i];}
\NormalTok{        xi[i].index = i + }\DecValTok{1}\NormalTok{;}
\NormalTok{    \}}
\NormalTok{    (}\DataTypeTok{void}\NormalTok{)qsort(xi, (}\DataTypeTok{size_t}\NormalTok{)nn, (}\DataTypeTok{size_t}\NormalTok{)}\KeywordTok{sizeof}\NormalTok{(}\KeywordTok{struct}\NormalTok{ couple), myComp);}
    \ControlFlowTok{for}\NormalTok{ (}\DataTypeTok{int}\NormalTok{ i = }\DecValTok{0}\NormalTok{; i < nn; i++) \{}
\NormalTok{        x[i] = xi[i].value;}
\NormalTok{        k[i] = xi[i].index;}
\NormalTok{    \}}
\NormalTok{    free(xi);}
\NormalTok{\}}
\end{Highlighting}
\end{Shaded}

\section{R code}\label{r-code}

\begin{Shaded}
\begin{Highlighting}[]
\KeywordTok{dyn.load}\NormalTok{(}\StringTok{"mySort.so"}\NormalTok{)}

\NormalTok{mySort <-}\StringTok{ }\ControlFlowTok{function}\NormalTok{ (x) \{}
\NormalTok{  h <-}
\StringTok{    }\KeywordTok{.C}\NormalTok{(}
      \StringTok{"mySort"}\NormalTok{,}
      \DataTypeTok{values =} \KeywordTok{as.double}\NormalTok{ (x),}
      \DataTypeTok{indices =} \KeywordTok{as.integer}\NormalTok{(}\DecValTok{1}\OperatorTok{:}\KeywordTok{length}\NormalTok{(x)),}
      \KeywordTok{as.integer}\NormalTok{(}\KeywordTok{length}\NormalTok{(x))}
\NormalTok{    )}
  \KeywordTok{return}\NormalTok{ (}\KeywordTok{list}\NormalTok{ (}\DataTypeTok{values =}\NormalTok{ h}\OperatorTok{$}\NormalTok{values, }\DataTypeTok{indices =}\NormalTok{ h}\OperatorTok{$}\NormalTok{indices))}
\NormalTok{\}}
\end{Highlighting}
\end{Shaded}


\end{document}
